\section{Architecture and Trust Boundaries}
\label{sec:04}
\label{sec:architecture}

\subsection{Architectural Overview}
\label{subsec:architecture-overview}

The architecture extends the previously developed biometric blockchain
voting system \citep{rifat2026biometric}. The inherited biometric
authentication layer identifies the voter. The private-voting extension
changes how an authenticated voter prepares and submits a ballot. Candidate
selection is separated from biometric authorization. Ballot verification
and tallying operate on encrypted data.
Table~\ref{tab:architecture-prior-comparison} identifies the architectural
changes relative to the previously published VoteChain system
\citep{rifat2026biometric}. The comparison concerns the respective voting
paths and their trust boundaries, not a controlled performance benchmark.
The legacy \texttt{LocalVoting} design and the present private-voting
extension perform different ballot-validation and tally operations.
The implementation comprises five connected parts. These are the biometric
authentication service, credential authorization and election management
services, the voter-facing web client with its privacy worker, the
Ethereum-compatible voting contracts, and the threshold trustees. Each part
has a defined role in the protocol. The web client prepares the encrypted
ballot. The voting contract checks its validity and maintains the encrypted
tally.

\begin{table}[tbp]
\centering
\caption{Voting-path architecture in the previously published VoteChain
system and the present VoteChain-ZK extension. The legacy column describes
the earlier published design; the current column describes the implemented
private-ballot pathway.}
\label{tab:architecture-prior-comparison}
\footnotesize
\setlength{\tabcolsep}{4pt}
\renewcommand{\arraystretch}{1.12}
\begin{tabular}{@{}>{\raggedright\arraybackslash}p{0.18\linewidth}>{\raggedright\arraybackslash}p{0.34\linewidth}>{\raggedright\arraybackslash}p{0.40\linewidth}@{}}
\toprule
\textbf{Aspect} & \textbf{Previous VoteChain (2026)} & \textbf{VoteChain-ZK (present study)} \\
\midrule
Biometric authorization &
Face/fingerprint verification authorizes voting through identity-aware services. &
Inherited face/fingerprint services authorize release of a protected voting credential; the circuit does not prove biometric matching. \\
\addlinespace[3pt]
Eligibility credential &
Salted pseudonymous voter hash, derived using the registered identity; checked against the contract's voter registry. &
Separate custom Poseidon credential commitment and frozen depth-17 Merkle root; the existing Semaphore registration anchor is not this proof root. \\
\addlinespace[3pt]
Candidate information flow &
The voting application prepares a transaction containing the selected candidate identifier. &
Candidate selection stays in the web client during biometric authorization and enters the local privacy worker for encryption. \\
\addlinespace[3pt]
Ballot submitted on-chain &
Candidate identifier and registered pseudonymous voter hash supplied to \texttt{LocalVoting.castVote}. &
Groth16 proof, election-scoped nullifier, and 16 encrypted candidate-slot ciphertexts supplied to \texttt{VoteChainDirect.castVote}; no plaintext candidate identifier. \\
\addlinespace[3pt]
Ballot acceptance &
Checks registered voter hash, unused vote status, and valid candidate identifier. &
Checks election phase, unused nullifier, canonical ciphertext coordinates, and a 73-public-signal Groth16 ballot-validity statement. \\
\addlinespace[3pt]
Tally during voting &
Plaintext \texttt{voteCount} incremented for the selected candidate. &
BabyJub exponential-ElGamal ciphertexts accumulated for fixed candidate slots; no plaintext per-candidate tally during voting. \\
\addlinespace[3pt]
Result disclosure &
Candidate counts directly readable from the contract; no trustee-decryption step. &
Three-of-five trustee partial decryptions checked off-chain; final totals and transcript hash recorded on-chain after separate state and arithmetic checks. \\
\addlinespace[3pt]
Residual trust &
Identity-aware service can reconstruct the mapping between voter identity and pseudonymous hash. &
Credential custodian can derive an identified voter's nullifier; ballot secrecy further depends on client integrity and trustee-share custody. \\
\bottomrule
\end{tabular}
\end{table}

Two commitment structures are present. A Semaphore identity and the
\texttt{ZKLocalVoting} contract support registration and election anchoring.
A separate custom Poseidon identity commitment supports the encrypted-voting
pathway. These commitments are derived from separately generated credentials
and are not interchangeable.
The private-voting workflow uses the custom commitment tree and
\texttt{VoteChainDirect}. Its eligibility proof is constructed against the
root configured for that contract. The \texttt{ZKLocalVoting} contract also
contains a candidate-indexed voting method. That method is separate from the
encrypted-ballot submission studied here.

% Source verification:
% votechain-central-server/controllers/voterController.js
% src/ZKLocalVoting.sol
% contracts/wp4-foundry/src/VoteChainWP4.sol

\subsection{Voter Registration and Election Configuration}
\label{subsec:architecture-registration}

Voter enrollment begins in the identity-aware registration service. The
service creates a Semaphore identity and a separate custom voting credential.
The custom credential contains a randomly generated secret. A Poseidon
identity commitment is derived from this secret. Both credentials are stored
in protected form alongside the enrolled voter record.
The two credentials serve different registration procedures. The Semaphore
commitment is associated with on-chain registration. The custom commitment
is retained for membership verification in the private-voting protocol. This
separation allows the encrypted-voting circuit to use its own eligibility
relation. The Semaphore group root is not used as the ballot-verification
root in \texttt{VoteChainDirect}.

Before election activation, the election management service collects the
registered custom commitments into an election-specific membership snapshot.
It constructs a fixed-depth Poseidon Merkle tree from these commitments. The
implementation uses a depth of 17. The resulting root identifies the
eligible commitment set for that election.
The service also prepares a candidate snapshot. Registered candidates are
ordered by candidate identifier and assigned fixed ballot slots. A
candidate-manifest hash binds this ordering to the election configuration.
The private ballot representation supports up to 16 candidate slots.

Trustee configuration is completed before election activation. Five trustees
participate in the configured distributed key-generation procedure. Its
public transcript contains verification material and the collective
encryption public key. The election management service verifies the
transcript and stores its canonical hash with the election record.
The election-start procedure deploys or resolves a Groth16 verifier and a
\texttt{VoteChainDirect} contract for the election. The election management
service configures the contract with the eligibility root, election
identifier, candidate manifest, candidate count, encryption public key, and
context-binding values.
The administrator then records the trustee-transcript hash on-chain while
the contract remains in its registration phase. After the configured values
are checked, the election is activated. Subsequent private ballots must
satisfy the frozen configuration.

% Source verification:
% votechain-central-server/services/zk/wp5PrivacyService.js
% votechain-central-server/services/elections/electionSnapshotService.js
% votechain-central-server/services/elections/electionTrusteeService.js
% votechain-central-server/services/elections/electionChainService.js

\subsection{Candidate-Blind Biometric Authorization}
\label{subsec:architecture-biometric-authorization}

The voter selects a candidate in the voter-facing web client. The selection
remains local during biometric authorization. The authorization request does
not transmit the candidate identifier.
The browser privacy worker first generates an ephemeral key pair for
credential release. Its public key is sent to the appropriate biometric
authorization endpoint. The private-voting pathway supports face and
fingerprint verification through separate routes. Both routes use the same
candidate-blind request structure.
The biometric authentication service processes the requested verification
through the corresponding acquisition and matching interface. After a
successful match, the credential authorization service identifies the
enrolled voter record. It checks whether the voter has an active custom
voting credential.

The credential authorization service retrieves the voter's membership path
from the frozen election snapshot. It reconstructs the Merkle root from the
registered commitment and path elements. This value must match the
eligibility root stored in \texttt{VoteChainDirect}. Authorization stops if
the values differ.
The credential authorization service then recovers the stored custom secret.
It protects the secret for delivery to the web client using ephemeral P-256
key agreement, HKDF-SHA-256, and AES-256-GCM. The authorization response
contains the encrypted credential envelope and Merkle witness. It also
identifies the election and specifies an authorization expiry time.

This separation has a specific privacy purpose. Biometric authentication
establishes which enrolled voter requests authorization. It does not require
the candidate selection. That selection is supplied to the privacy worker
after authorization, when the encrypted ballot is constructed.
The biometric and credential services remain identity-aware components. The
credential authorization service can associate a registered voter with the
stored credential and recover its secret during authorization. The resulting
identity-linkage boundary is defined in Section~\ref{sec:system-model}.

% Source verification:
% votechain-central-server/middleware/wp5AuthorizationSecurity.js
% votechain-central-server/controllers/wp5Controller.js
% votechain-central-server/controllers/fingerprintController.js
% votechain-central-server/services/zk/wp5PrivacyService.js
% votechain-frontend-ts/src/components/Biometric/VerificationModal.tsx

\subsection{Local Encrypted Ballot Construction and Submission}
\label{subsec:architecture-ballot-submission}

After authorization, the web client passes the credential envelope and
membership witness to its privacy worker. The worker uses the ephemeral
private key to recover the custom voting secret. It reconstructs the
identity commitment and verifies the supplied Merkle path against the
configured eligibility root.
The worker checks the election identifier and candidate manifest. It also
recomputes the election-context binding using the configured candidate
count, contract domain, and encryption public key. These checks stop ballot
preparation if the public election data are inconsistent.

The selected candidate is represented as a one-hot vector. Exactly one
active candidate slot contains the value one. The remaining slots contain
zero. The vector retains the fixed 16-slot representation used by the
voting circuit.
The worker encrypts each slot using exponential ElGamal over BabyJub. It
generates fresh encryption randomness for the individual slots. Each
ciphertext contains two curve points. The candidate selection and encryption
randomness remain private proof inputs.

The worker computes the election-scoped nullifier from the credential secret
and election identifier. It then generates a Groth16 proof using the custom
Circom circuit. The proof binds eligibility, election context, ballot
structure, and ciphertext construction to the submitted public values
\citep{groth2016snark}.
The worker checks the returned 73 public signals against the expected
election values, nullifier, and ciphertext coordinates. It then verifies
the generated proof locally. Any mismatch or failed proof verification
stops submission.

The web client passes the verified proof, nullifier, and ciphertext
coordinates to the wallet transaction interface. The transaction invokes
\texttt{VoteChainDirect.castVote}. It does not include a plaintext candidate
identifier.
The contract checks that the election is active and the nullifier has not
been used. It also checks the canonical representation of the ciphertext
coordinates. The Groth16 verifier evaluates the proof against the
contract's configured public inputs.
Once these checks succeed, the contract records the nullifier and adds the
ciphertext to the encrypted tally. It increments the accepted-ballot count
and emits an acceptance event. The event records the nullifier, a ciphertext
hash, and the updated ballot count rather than the selected candidate.

% Source verification:
% votechain-frontend-ts/src/components/Voter/VoteCasting.tsx
% votechain-frontend-ts/src/privacy/wp5Privacy.worker.ts
% votechain-frontend-ts/src/services/wp5Direct.ts
% circuits/VoteChainVoteDirect.circom
% circuits/lib/VoteChainVoteCore.circom
% contracts/wp4-foundry/src/VoteChainWP4.sol

\subsection{Encrypted Tallying and Result Finalization}
\label{subsec:architecture-finalization}

The voting contract maintains an encrypted aggregate for each of the 16
candidate slots. After an accepted ballot, the corresponding ciphertext
points are added to the current aggregate using BabyJub curve operations.
The accepted-ballot count is updated separately.
During active voting, the contract stores encrypted aggregates rather than
plaintext candidate totals. This encrypted state is publicly inspectable.
Recovering the candidate counts requires the corresponding decryption
capability.

After voting ends, the administrator changes the contract state from Active
to Ended. The election management service reads the encrypted aggregates and
accepted-ballot count from \texttt{VoteChainDirect}. These values are used
to check the threshold-finalization transcript.
The finalization protocol uses a three-out-of-five trustee threshold.
Trustees produce partial decryptions and accompanying verification records
for the encrypted candidate aggregates. An external verifier checks the
distributed key-generation transcript and the partial-decryption evidence.
For each candidate slot, the verifier requires at least three valid trustee
contributions. It combines those contributions and recovers the aggregate
count within the permitted vote-count range. The election management service
also checks that the finalization transcript matches the encrypted tally and
accepted-ballot count recorded on-chain.

After successful external verification, the administrator submits the
recovered candidate totals and finalization-transcript hash to
\texttt{VoteChainDirect}. The contract requires the election to be in its
Ended phase and a trustee transcript to have been pinned beforehand.
The contract enforces additional checks. Candidate slots outside the active
manifest must have zero totals. No candidate total may exceed the
accepted-ballot count. The sum of the active candidate totals must equal the
number of accepted ballots.
If these conditions hold, the contract records the final totals and
transcript hash. It then changes the election state to Finalized. Plaintext
candidate results become accessible through its finalized-result interface.

The contract does not directly verify trustee partial-decryption proofs.
Those checks are performed by the external transcript verifier before the
finalization transaction is submitted. The contract's arithmetic checks
alone therefore do not prove that the submitted totals were derived from
valid threshold decryptions.
Independent verification requires the complete transcript and published
encrypted aggregate. An auditor can check the trustee evidence and compare
the verified transcript hash with the on-chain commitment. This provides an
externally auditable result-disclosure procedure under the stated trust and
transcript-availability conditions.
The architecture connects biometric authorization, private ballot
construction, proof-verified encrypted accumulation, and threshold-based
result disclosure. The following section specifies the cryptographic
relations and verification operations used at these stages.

% Source verification:
% crypto/src/wp6-threshold.mjs
% votechain-central-server/services/elections/electionFinalizationService.js
% contracts/wp4-foundry/src/VoteChainWP4.sol

